\documentclass{article}
\usepackage{amsmath}
\usepackage{algorithm}
\usepackage{algorithmic}
\usepackage[utf8]{inputenc}
\usepackage[spanish]{babel}

\title{Un algoritmo reutilizable: m\'etodo de la secante}
\date{}

\begin{document}
\maketitle

Se aproximan las ra\'ices de dos funciones con el mismo algoritmo.
Como en Newton--Raphson, se trata de un m\'etodo iterativo, pero la
secante aproxima la derivada a partir de dos puntos.

\begin{equation}
    f(x) = x^2 - 2
    \label{eq:calc:f}
\end{equation}

\begin{equation}
    g(x) = x^3 - 2
    \label{eq:calc:g}
\end{equation}

\begin{algorithm}
\caption{M\'etodo de la secante}
\label{alg:calc:secante}
\begin{algorithmic}
    \REQUIRE $F, a, b, tol, limite$
    \STATE $i = 0$
    \WHILE{$i < limite \AND |F(b)| > tol$}
        \STATE $s = b - F(b) \cdot (b - a) / (F(b) - F(a))$
        \STATE $a = b$
        \STATE $b = s$
        \STATE $i = i + 1$
    \ENDWHILE
    \RETURN $b$
\end{algorithmic}
\end{algorithm}

La etiqueta identifica la funci\'on \texttt{secante}; el algoritmo se
ejecuta al llamarla, no al declararla. Sus par\'ametros son los indicados
en \texttt{Require}, en el mismo orden.

\begin{equation}
    raizuno = secante(f, 1, 2, 1e-8, 50)
    \label{eq:calc:raizuno}
\end{equation}

\begin{equation}
    raizdos = secante(g, 1, 2, 1e-8, 50)
    \label{eq:calc:raizdos}
\end{equation}

Las ra\'ices aproximadas son:

\begin{equation}
    raizuno =
    1.414214
    \label{eq:res:raizuno}
\end{equation}

\begin{equation}
    raizdos =
    1.259921
    \label{eq:res:raizdos}
\end{equation}

Si se alcanza el l\'imite de iteraciones, el algoritmo devuelve la
\'ultima aproximaci\'on sin garantizar la convergencia. La f\'ormula
requiere que $F(b) - F(a)$ no sea cero.

\end{document}
